\chapter{Mobile Application}

The mobile application, built with Flutter for cross-platform deployment,
performs three roles: it manages the BLE link to the glove, it runs the
on-device neural network, and it presents recognised gestures as text and
speech.

\section{Architecture}

The application is organised into services with clear responsibilities: a BLE
service that handles scanning, connection, and packet parsing; an inference
service that buffers sensor frames and runs the model; a text-to-speech
service; and a preferences service that holds the per-model configuration. The
main screen is a single capture-and-display view; a segmented control switches
between the \emph{Letters} and \emph{Words} models.

\section{Bluetooth Link and Packet Parsing}

On connection the BLE service requests the appropriate runtime permissions,
scans for the glove by address (falling back to its advertised name), connects,
and—crucially—negotiates a larger MTU before subscribing to the notify
characteristic (Chapter~4). Each 180-byte notification is parsed into 45
little-endian floats forming one sensor frame, which is pushed into a fixed-%
size rolling buffer. Frames are rejected if their length is wrong or any value
is non-finite. The Android manifest declares the modern \texttt{BLUETOOTH\_%
SCAN} and \texttt{BLUETOOTH\_CONNECT} permissions; an early bug, in which the
app required a location permission that the manifest no longer declared on
recent Android, was traced to this and fixed.

\section{Gated Capture}

Rather than classifying continuously, the application uses \emph{gated
capture}, mirroring the data-collection procedure exactly. The user presses a
\textbf{Listen} button; the screen shows a brief countdown so the signer can
prepare, then a clear ``sign now'' cue while a fresh window of sensor frames is
collected; the model is then run once on that window. This guarantees that the
window seen at inference matches the window used during training, and it makes
the device's behaviour predictable and demonstrable. It also matches how
fingerspelling is naturally performed: one deliberate gesture at a time.

\section{On-Device Inference Pipeline}

When a window is captured, the inference service applies, in order: (i) the
per-model feature transform (the relative transform for a relative model, or
none for an absolute model), (ii) standard-scaler normalisation using the
model's exported statistics, (iii) quantisation of the normalised values to
int8 using the input tensor's scale and zero-point, (iv) execution of the
TensorFlow Lite model, and (v) dequantisation of the int8 output back to class
probabilities, followed by an arg-max. Each of these steps must agree exactly
with the training pipeline; the feature transform and the normalisation
statistics are therefore selected \emph{per model} (Section~\ref{sec:permodel}).

\section{Per-Model Configuration}\label{sec:permodel}

The Letters and Words models have different class sets, different normalisation
statistics, and—in principle—could use different feature transforms. The
application therefore swaps the model file, the label list, the normalisation
statistics, and the transform flag together whenever the user changes mode. A
subtle bug class this avoids is silently using one model's normalisation with
another model's weights; the normaliser is reset on each load so a failed swap
can never leave stale statistics in place.

\section{Output}

In Words mode, a recognised sign is spoken immediately through text-to-speech.
In Letters mode, recognised letters accumulate into a text buffer so the user
can fingerspell a word; a \textbf{Speak} button reads the assembled word aloud
and a \textbf{Clear} button empties the buffer to start again. A confidence
threshold suppresses low-confidence predictions so that an uncertain capture
asks the user to try again rather than committing a wrong letter.

Figure~\ref{fig:app} shows the application in use. In Words mode (left) the
recognised sign is shown in large type with a confidence bar beneath it and is
spoken at once; multi-word signs such as \textsc{how are you} are displayed and
spoken as spaced phrases rather than the run-together tokens the model emits. In
Letters mode (centre) each accepted letter is appended to the buffer shown at the
top, and the \textbf{Speak} and \textbf{Clear} controls let the user read back
or reset the assembled word. The Settings screen (right) exposes the confidence
threshold, the text-to-speech voice, and the Letters-mode auto-clear interval.
The same confidence threshold governs both modes, so the interface stays quiet
rather than guessing when a capture is ambiguous.

% FIGURE — mobile app screenshots (Words, Letters, Settings).
\begin{figure}[!htbp]
    \centering
    \includegraphics[width=0.3\textwidth]{AppSS2.png}
    \hfill
    \includegraphics[width=0.3\textwidth]{AppSS1.png}
    \hfill
    \includegraphics[width=0.3\textwidth]{AppSS3.png}
    \caption{\label{fig:app}The mobile application. Left: Words mode, which speaks
    a recognised sign immediately. Centre: Letters mode, which accumulates
    fingerspelled letters into a buffer with \textbf{Speak} and \textbf{Clear}
    controls. Right: the Settings screen, exposing the confidence threshold,
    text-to-speech voice, and the Letters-mode auto-clear interval.}
\end{figure}
